\begin{pastexam}[label={px:nccu-114-dsa-04}]{政大資科所 114 學年度 第 4 題}
（單選；每題 4 分）

(4\%) Which of the following statement concerning hashing is false?
\begin{enumerate}[label=\Alph*.]
\item While hash function is designed based on arithmetic operation, hashing is also suitable for
strings.
\item If there is no collision, the worst time complexity for searching an element in a hash table
is O(1).
\item A good hash function should produce the same hash value for similar keys.
\item Hashing is typically used for symbol tables in compilers.
\item The size of the hash table should ideally be a prime number to reduce the likelihood of
collisions.
\end{enumerate}
\end{pastexam}

\begin{pastsoln}
答案是 C。

逐項檢查：
\begin{itemize}
\item A.~對。 字串先轉成整數再雜湊，例如把字元當成某個底數的各位數字，用 Horner 法則算
$\sum c_i \cdot b^{k-1-i} \bmod m$（資結書第~12.1~節）。 雜湊函數是算術運算，但輸入可以是字串。
\item B.~對。 沒有碰撞時，每個鍵算出的位置裡只有它自己，搜尋就是算一次雜湊值、看一格，
$O(1)$（假設算雜湊值本身是 $O(1)$）。
\item C.~錯，這就是答案。 好的雜湊函數要讓鍵\textbf{均勻分散}到各個位置（簡單均勻雜湊的假設，資結書
第~12.1~節）。 實際的鍵常常很像，例如變數名稱 \texttt{x1}、\texttt{x2}、\texttt{x3}，如果相似的鍵
算出相同的雜湊值，它們全部擠在同一格，搜尋退化成線性。 好的雜湊函數應該讓相似的鍵得到\textbf{不同}的
雜湊值。
\item D.~對。 編譯器的符號表要用名稱快速查出變數的型別、位址等資料，插入和查詢都很頻繁、幾乎
不需要依序走訪，是雜湊表的典型用途。
\item E.~對。 除法雜湊 $h(k) = k \bmod m$ 時，$m$ 取質數（而且不要太靠近 $2$ 的次方），鍵的規律
（例如都是偶數、都是某個數的倍數）比較不會讓它們集中在少數幾格。 這是教科書的標準建議。
\end{itemize}
\end{pastsoln}
